% ============================================================
%  CHAPTER 2 – System Analysis
% ============================================================
\chapter{System Analysis}
\label{ch:system-analysis}
This chapter describes the background of the design~\ref{sec:background}, provides the problem statement~\ref{sec:problem-statement}, discusses project objectives~\ref{sec:objectives}, as well as the system requirements~\ref{sec:requirements}.

\section{Background}
\label{sec:background}

[Describe the domain context and motivation. Explain why the problem
exists, what current approaches are insufficient, and why an engineering
solution is warranted. Cite relevant sources.]

\section{Problem Statement}
\label{sec:problem-statement}

[Clearly state the engineering problem to be solved in one or two
focused paragraphs. Avoid proposing solutions here.]

\section{Project Objectives}
\label{sec:objectives}

The objectives of this project are:
\begin{itemize}
    \item \textbf{Objective 1:} \ldots
    \item \textbf{Objective 2:} \ldots
    \item \textbf{Objective 3:} \ldots
\end{itemize}

\section{Requirements}
\label{sec:requirements}

The system requirements are based on the document
\textit{Engineering Design 2026 -- Hoisting System System
Specification}, Revision 0, dated 4 August 2026. The specification
defines the functional, performance, interface, physical, design,
environmental, and other quality requirements for the
rail-mounted gantry/trolley hoisting system.

The system comprises five work packages: MCTR I (Lateral
Movement), MCTR II (Spooling), EEII I (Drive Electronics),
EEII II (Magnetic Load Attachment), and REII (Remote Control
Dashboard). The work packages interface with a shared embedded
controller based on the ESP32-S3-WROOM-1 module and an isolated
I/O board.

The requirements below summarise the applicable system
specification. Detailed implementation requirements and
controller pin assignments are defined in the relevant work-package
scopes of work and Controller Specification.

\subsection{Functional Requirements}
\label{subsec:func-req}

The functional requirements define the operations that the system
and its constituent subsystems shall perform.

\begin{itemize}
    \item \textbf{FR1 -- Lateral movement:}
    Upon receipt of a motion command from the controller, the
    lateral movement mechanism shall move the trolley along the
    I-beam rail in the commanded direction using a commercial
    off-the-shelf (COTS) DC motor costing less than R250.
    The associated performance target is defined by PT1.

    \item \textbf{FR2 -- Position feedback:}
    The lateral movement mechanism shall report wheel position
    to the controller via quadrature encoder feedback to support
    motion profiling by REII.

    \item \textbf{FR3 -- Wheel manufacturing:}
    The lateral movement wheels shall be custom-manufactured
    by 3D printing to support the wheel wear-characterisation
    exercise.

    \item \textbf{FR4 -- Hoisting:}
    Upon receipt of a raise or lower command, the spooling
    mechanism shall drive the hoist gearbox to lift or lower
    a payload of up to 20\,kg in the commanded direction.
    The associated performance target is defined by PT2.

    \item \textbf{FR5 -- Hoist load feedback:}
    The spooling mechanism shall continuously report the
    on-hook load to the controller during operation.

    \item \textbf{FR6 -- Fail-safe braking:}
    The active brake shall engage automatically upon loss
    of drive power or command signal and shall hold the
    payload without drift.

    \item \textbf{FR7 -- Motor drive:}
    Upon receipt of gate/PWM commands from the controller,
    the drive electronics shall convert single-phase
    230\,VAC mains power to a DC bus and invert it to a
    three-phase drive for the induction motor.

    \item \textbf{FR8 -- Drive protection:}
    The drive electronics shall detect and interrupt drive
    output during shoot-through or overcurrent conditions
    to protect the switching devices from damage.

    \item \textbf{FR9 -- Drive current feedback:}
    The drive electronics shall provide a single-ended
    current-sense signal, capped below 2.048\,V, to the
    controller for firmware-based fault inference.

    \item \textbf{FR10 -- Magnetic load attachment:}
    Upon receipt of an engage command, the magnetic load
    attachment shall generate sufficient holding force
    through a three-point electromagnet to retain a payload
    of up to 20\,kg under static conditions. The associated
    performance target is defined by PT3.

    \item \textbf{FR11 -- Magnetic load indication:}
    The magnetic load attachment shall report an on-load
    indication to the controller using current-sense
    discrimination between a loaded and unloaded grip.

    \item \textbf{FR12 -- Dashboard telemetry:}
    Upon receipt of telemetry from the controller via the
    nRF24L01 link, the dashboard shall display all operational
    parameters to the operator, including position, speed,
    load, current, fault status, and grip status.

    \item \textbf{FR13 -- Dashboard control:}
    The dashboard shall provide the operator with control
    of all system operations, including lateral motion,
    hoisting, and magnet-engagement commands.

    \item \textbf{FR14 -- Independent real-time control:}
    The embedded controller shall execute its real-time
    control loop, including drive-signal generation, sensor
    acquisition, and safety interlocks, independently of
    wireless-link or dashboard availability.

    \item \textbf{FR15 -- Standalone operation:}
    For groups without computer-engineering capacity, the
    controller shall accept motion and control commands
    via UDP over its built-in Wi-Fi interface without
    requiring additional hardware.
\end{itemize}


\subsection{Performance Requirements}
\label{subsec:perf-req}

The specification defines three quantitative performance targets.
These targets are used for graduated performance assessment rather
than as binary functional pass/fail requirements.

\begin{itemize}
    \item \textbf{PT1 -- Lateral movement speed:}
    The trolley shall achieve a lateral speed of at least
    1\,m/s, measured at a distance of 1\,m from the departure
    point under unloaded conditions. Performance marks scale
    with the demonstrated speed relative to the target.

    \item \textbf{PT2 -- Hoisting performance:}
    The hoisting mechanism shall complete a 1.6\,m hoist
    stroke within 20\,s while carrying a 20\,kg payload,
    corresponding to an average hoisting rate of approximately
    80\,mm/s. Performance marks scale with the demonstrated
    hoisting rate.

    \item \textbf{PT3 -- Magnetic holding performance:}
    The magnetic load attachment shall retain a 20\,kg payload
    against a specified disturbance force. The numerical
    disturbance-force threshold remains to be determined
    through further dynamics modelling. Performance marks
    scale with the demonstrated holding performance relative
    to the eventual threshold.
\end{itemize}


\subsection{Physical Requirements}
\label{subsec:phys-req}

The physical requirements define the payload interface and
essential hardware safety provisions.

\begin{itemize}
    \item \textbf{PHY1 -- Payload interface:}
    The payload used for testing and demonstration shall have
    a flat sheet-metal attachment surface measuring
    600\,mm $\times$ 400\,mm and at least 3\,mm thick for
    engagement by the magnetic load attachment. Additional
    mass may be used to increase the total payload to a
    maximum of 20\,kg, provided the combined centre of mass
    remains centred on the attachment surface. This
    requirement does not require the general system mounting
    plate to be made from sheet metal.

    \item \textbf{PHY2 -- Hardware emergency stop:}
    The system shall include a hardware emergency-stop
    mechanism wired directly into the high-voltage (HV) bus.
    It shall be capable of removing power from the drive
    electronics independently of the software or firmware
    state.

    \item \textbf{PHY3 -- Mains protection:}
    The mains supply feeding the drive electronics shall be
    protected by a residual-current device (RCD), considering
    the combination of high-voltage wiring and outdoor
    exposure to moisture.
\end{itemize}


\subsection{Non-Functional Requirements}
\label{subsec:nonfunc-req}

In the system specification, these requirements are classified
under \textit{Other Qualities}. They address reliability,
maintainability, usability, and affordability.

\begin{itemize}
    \item \textbf{OQ1 -- Reliability:}
    Shoot-through and overcurrent protection shall be
    validated on a milled test board before construction
    of the full-scale variable-frequency drive (VFD)
    assembly.

    \item \textbf{OQ2 -- Maintainability:}
    The core controller board shall be designed as reusable,
    general-purpose laboratory infrastructure, providing
    local analogue acquisition and data logging independently
    of any single year's design problem.

    \item \textbf{OQ3 -- Usability:}
    The dashboard shall present telemetry and controls so
    that an operator can safely command the system after
    a brief orientation, without requiring specialist RF
    or embedded-systems knowledge.

    \item \textbf{OQ4 -- Affordability:}
    Component costs shall remain within the stated constraints
    for each work package, including a cost below R250 for
    the lateral-movement motor. The total system cost per
    group shall not exceed R5,000, inclusive of all
    work-package procurement.
\end{itemize}


\subsection{Design and Construction Requirements}
\label{subsec:design-req}

The following design requirements constrain the permitted
implementation choices.

\begin{itemize}
    \item \textbf{DR1 -- Lateral-movement motor:}
    The lateral-movement drive motor shall be a COTS DC motor
    costing less than R250.

    \item \textbf{DR2 -- Spooling gearbox:}
    The spooling gearbox shall not use a worm-gear configuration.
    The spooling work-package scope further specifies a
    minimum two-stage gearbox.

    \item \textbf{DR3 -- Electromagnet construction:}
    The magnetic load attachment shall not use a complete
    COTS electromagnet or solenoid. COTS core material is
    permitted, but the electromagnets must be manufactured
    as part of the project.

    \item \textbf{DR4 -- Wireless module:}
    The wireless link shall use the nRF24L01 module.

    \item \textbf{DR5 -- Embedded controller:}
    The embedded controller shall be based on the
    ESP32-S3-WROOM-1 (N8R2) module.
\end{itemize}


\subsection{Interface Requirements}
\label{subsec:interface-req}

The controller shall provide the interfaces required to
coordinate the five work packages.

\begin{itemize}
    \item \textbf{IR1 -- Drive electronics interface:}
    The controller shall provide an isolated six-line
    gate/PWM interface to EEII I for three-phase motor
    control, together with one reserve line. A current-sense
    input shall also be provided.

    \item \textbf{IR2 -- Magnetic attachment interface:}
    The controller shall provide an isolated PWM interface
    and a single-ended current-sense input of no more than
    2.048\,V to EEII II.

    \item \textbf{IR3 -- Lateral movement interface:}
    The controller shall provide direction, PWM, and
    quadrature-encoder interfaces to MCTR I.

    \item \textbf{IR4 -- Spooling interface:}
    The controller shall provide a load-cell input and
    brake-control interface to MCTR II. The load-cell
    acquisition path remains to be finalised between an
    ADS1220 differential input and a dedicated HX711.

    \item \textbf{IR5 -- Expansion interface:}
    The controller shall provide a five-pin, protocol-agnostic
    expansion header for group-specific extensions. The
    header shall be documented as a low-voltage digital-logic
    interface only.
\end{itemize}


\subsection{Environmental and Resource Requirements}
\label{subsec:environmental-req}

The system shall satisfy the following environmental and
external-resource requirements.

\begin{itemize}
    \item The system shall operate on the existing I-beam rail
    and gantry structure located outdoors on the NWU North Campus.

    \item The system shall operate outdoors under ambient
    weather conditions typical of the site, including
    temperature variation, precipitation, solar and UV
    exposure, wind, and dust. It shall also operate in the
    presence of other groups' systems on shared or adjacent
    rail infrastructure.

    \item Electrical enclosures and wiring shall provide
    protection against rain, dust, and incidental moisture.
    Exposed materials, including 3D-printed or polymer
    components, shall be selected or rated for sustained
    outdoor UV and temperature exposure over the semester.

    \item The system shall remain operational throughout the
    design, build, integration, and demonstration phases of
    the 2026 academic semester.

    \item The system shall receive electrical power through
    one dedicated single-phase, 50\,Hz, 230\,VAC mains
    supply point per group.

    \item System-level integration, including the main
    on/off switch and distribution of mains power to the
    HV bus and to the selected COTS low-voltage supplies
    of the individual work packages, shall be assigned
    to EEII I.
\end{itemize}


\subsection{Operating States and Modes}
\label{subsec:states-modes}

The system specification defines three operating states and
two operating modes.

\textbf{Operating states:}
\begin{itemize}
    \item \textbf{Idle:} The system is powered but is not
    executing a commanded motion, hoisting, or
    magnet-engagement action.

    \item \textbf{Operational:} The system is actively
    executing a commanded lateral movement, hoisting, or
    magnet-engagement function under either operating mode.

    \item \textbf{Fault / E-Stopped:} The hardware emergency
    stop has removed power from the HV bus. All motive
    functions are inhibited independently of the software
    or firmware state.
\end{itemize}

\textbf{Operating modes:}
\begin{itemize}
    \item \textbf{Dashboard-controlled:} The operator controls
    the system through the REII PC-based dashboard over the
    nRF24L01 wireless link.

    \item \textbf{Standalone:} The operator controls the system
    directly through the controller's built-in Wi-Fi/UDP
    interface, without requiring a dashboard.
\end{itemize}

At power-up, the system shall begin in the Idle state. Both
operating modes shall be available in the Operational state.
The Fault / E-Stopped state shall take precedence over either
mode and shall inhibit motive outputs regardless of the mode
active when the emergency stop was activated.


\subsection{Open Requirements and Verification}
\label{subsec:open-req-verification}

The system specification identifies the following items as
requiring further definition or implementation-level verification.

\begin{itemize}
    \item The numerical disturbance-force threshold for PT3
    remains open pending further dynamics modelling.

    \item The spooling load-cell acquisition path remains an
    implementation choice between an ADS1220 differential
    input and a dedicated HX711.

    \item The system requirements must be assigned explicit
    verification methods, including defined test procedures,
    measurement methods, and acceptance criteria.

    \item Shoot-through and overcurrent protection shall be
    validated on the milled test board before the full-scale
    VFD assembly is constructed.
\end{itemize}

Functional requirements shall be verified as pass/fail
against their stated behaviour. Performance targets PT1--PT3
shall be assessed according to the degree to which the
specified targets are achieved. Open numerical thresholds
shall not be treated as fully defined acceptance criteria
until they have been resolved.


\section{Risk Analysis}
\label{sec:risk-analysis}

The risk analysis shall identify hazards and implementation
risks associated with the mechanical, electrical, magnetic,
and control subsystems. Particular attention shall be given
to the high-voltage drive electronics, outdoor operation,
payload retention, moving mechanisms, and emergency-stop
functionality.

The risk register shall document each identified risk,
its likelihood and impact, mitigation measures, and any
residual risk following mitigation. The detailed risk
register is provided in Appendix~\ref{app:risk}.

